% =============================================================================
%  Notas del Profesor — Sesión 03: Sistemas de Numeración
%  Programación Avanzada — Universidad Panamericana
%  Compilar: pdflatex sesion03_sistemas_numeracion_profesor.tex
% =============================================================================
\documentclass[12pt, oneside]{article}
\usepackage[profesor]{progavanzada}   % versión profesor (portada azul UP)

\begin{document}

\portada{Sesión 03}{Sistemas de Numeración}{2026}

\tableofcontents
\newpage

% =============================================================================
\section{Objetivos de la sesión}
% =============================================================================

Al finalizar la sesión el alumno será capaz de:

\begin{enumerate}
  \item Explicar por qué distintas culturas adoptaron distintas bases de
        numeración, con al menos tres ejemplos históricos.
  \item Convertir números entre base 2, 8, 10 y 16 usando los métodos
        de suma de pesos y divisiones sucesivas.
  \item Justificar por qué las computadoras operan en base 2 y por qué
        el hexadecimal facilita la lectura del binario.
  \item Resolver ejercicios de conversión directa entre binario, octal
        y hexadecimal sin pasar por decimal.
\end{enumerate}

\begin{notaprofesor}[Duración estimada]
  50 minutos. Sugerido: 15 min de apertura histórica (conversación
  + preguntas) + 25 min de desarrollo técnico en cuaderno + 10 min de
  práctica guiada con los ejercicios.
\end{notaprofesor}

% =============================================================================
\section{Secuencia de clase}
% =============================================================================

\subsection{Apertura — El por qué de los sistemas de numeración (15 min)}

Abrir \textbf{sin escribir nada} en el pizarrón. Lanzar la pregunta al grupo:

\begin{notaprofesor}[Pregunta de arranque]
  \textit{``¿Por qué usamos exactamente diez dígitos y no nueve o doce?''}\\[4pt]
  Deja que la conversación fluya. La respuesta que buscas es ``por los dedos'',
  pero no la des tú primero. Si el grupo no llega, guíalos: ``¿qué tienen
  en común todos los humanos que aprendieron a contar hace miles de años?''
\end{notaprofesor}

Una vez establecido el origen anatómico de la base 10, introduce los tres
ejemplos históricos en orden cronológico:

\begin{enumerate}
  \item \textbf{Mayas — base 20.}
        Pregunta: ``¿qué contaban los mayas que los humanos modernos no
        cuentan?'' (dedos de los pies). Menciona el cero maya: fue una de
        las primeras civilizaciones en usarlo como valor posicional, siglos
        antes que Europa.

  \item \textbf{Babilonios — base 60.}
        Conecta inmediatamente con algo que los alumnos ya conocen: reloj,
        minutos, segundos, grados en un círculo. La pregunta retórica que
        funciona muy bien: ``¿alguien sabe por qué un minuto tiene 60 segundos
        y no 100?''. La respuesta es que lo heredamos de los babilonios, y
        ellos lo eligieron porque 60 tiene 12 divisores enteros.

  \item \textbf{Propuesta de base 12.}
        Menciona que durante la Revolución Francesa, mientras se diseñaba el
        sistema métrico, varios matemáticos argumentaron en favor de base 12
        por su divisibilidad. Ganar el argumento de ``más fracciones exactas''
        es fácil con la tabla de divisores.
        La Dozenal Society (\url{https://www.dozenalsociety.org.uk}) sigue
        activa hoy en día.
        Pregunta de cierre: ``¿Con qué ya contamos en docenas en el día a día?''
        (meses, horas, huevos).
\end{enumerate}

\begin{notaprofesor}[Tono de esta parte]
  El objetivo no es que memoricen fechas: es que entiendan que la base
  no es algo ``natural'' ni único, sino una \textbf{elección conveniente}.
  Eso abre la puerta a explicar por qué las computadoras eligieron base 2,
  que es la pregunta que deberían hacerse solos antes de que tú la respondas.
\end{notaprofesor}

\subsection{Transición — ¿Qué elegiría una computadora? (5 min)}

Antes de explicar el binario, lanza la pregunta:

\begin{notaprofesor}[Pregunta de transición]
  \textit{``Si una computadora contara en base 10 igual que nosotros,
  ¿cómo distinguiría el número 3 del número 4 usando electricidad?''}\\[4pt]
  Respuesta que buscas: necesitaría diez niveles de voltaje distintos
  (0V, 1V, \ldots, 9V), lo cual es muy propenso a errores por ruido
  eléctrico. Con solo dos niveles (bajo / alto), la distinción es robusta.
  De ahí el binario.
\end{notaprofesor}

\subsection{Desarrollo técnico — Binario, octal y hexadecimal (25 min)}

Trabaja sobre el cuaderno o en el pizarrón. El orden recomendado:

\begin{enumerate}
  \item Sistema posicional genérico: peso de cada posición como
        $d_i \times b^i$. Un ejemplo en base 10 primero para anclar
        el concepto.
  \item \textbf{Binario → decimal}: sumas de potencias de 2.
        Haz uno junto con el grupo; deja que ellos hagan el siguiente.
  \item \textbf{Decimal → binario}: divisiones sucesivas.
        Muestra el proceso completo con $25_{10}$ o $1521_{10}$.
  \item \textbf{Octal}: igual que binario pero con base 8. Énfasis en la
        conversión directa octal ↔ binario por grupos de 3 bits.
  \item \textbf{Hexadecimal}: introduce los símbolos A–F. Énfasis en la
        conversión directa hex ↔ binario por grupos de 4 bits.
\end{enumerate}

\begin{notaprofesor}[Por qué octal antes que hexadecimal]
  Octal es más fácil de interiorizar (los dígitos son 0–7, sin letras).
  Sirve como puente: ya que los alumnos entienden la conversión directa
  con grupos de 3 bits, la de hex con grupos de 4 bits es inmediata.
\end{notaprofesor}

\subsection{Ejercicios (10 min)}

Los seis ejercicios de conversión al final de las notas del alumno están
ordenados de menor a mayor dificultad.

\begin{enumerate}
  \item Decimal → binario (el algoritmo de divisiones).
  \item Binario → decimal (suma de pesos, número largo).
  \item Binario → hex (agrupación por 4 bits).
  \item Hex → binario (sustitución directa).
  \item Octal → hex (vía binario).
  \item Base 4 → hex (vía binario).
\end{enumerate}

Deja el problema de concurso (Carina y la calculadora) como tarea o para
los alumnos que terminen rápido. Es un problema de Olimpiada que requiere
darse cuenta de que contar sin 2 ni 7 equivale a contar en base 8.

\begin{notaprofesor}[Solución del problema de Carina]
  Los dígitos disponibles son $\{0, 1, 3, 4, 5, 6, 8, 9\}$: exactamente 8
  símbolos. Si los renombramos $\{0,1,2,3,4,5,6,7\}$, Carina en realidad
  está contando en base 8. El milésimo número en base 8 es $1750_8$,
  que en los dígitos de Carina es \textbf{1, 9, 5, 0} $\to$ \textbf{1950}.
  (Verificar: $1\times 8^3 + 7\times 8^2 + 5\times 8 + 0 = 512+448+40 = 1000$.)
\end{notaprofesor}

% =============================================================================
\section{Contexto histórico — Notas adicionales}
% =============================================================================

Esta sección amplía los datos históricos para responder preguntas del grupo
o enriquecer la explicación.

\subsection{Mayas}

El sistema maya era \textbf{vigesimal y posicional}, con una modificación:
en el conteo del calendario, la segunda posición valía 18 × 20 = 360 (para
aproximar el año solar) en lugar de $20^2 = 400$. Esto lo hace un sistema
mixto en el contexto calendárico, aunque en el contexto aritmético puro
era base 20 consistente.

El símbolo del cero maya (una concha o caracol) data aproximadamente del
siglo \textsc{iv} d.C.\ en los registros escritos conocidos, aunque podría
ser más antiguo en la tradición oral.

\begin{notaprofesor}[Referencia]
  Ifrah, G. \textit{Historia Universal de las Cifras}. Espasa Calpe, 1994.
  Caps. 22–23 (numeración maya). Disponible en varias bibliotecas
  universitarias; es el texto de referencia estándar en historia de las
  matemáticas.
\end{notaprofesor}

\subsection{Babilonios}

Los sumerios (c.\ 3000 a.C.) desarrollaron escritura cuneiforme y, con ella,
un sistema de conteo en base 60. Existen tablillas de arcilla con tablas de
multiplicación en base 60 de hace más de 4\,000 años.

Una teoría complementaria sobre la elección de 60: los sumerios pueden haber
combinado un sistema de base 5 (dedos de una mano) con uno de base 12
(contando con el pulgar las 3 falanges de cada uno de los otros 4 dedos:
$4 \times 3 = 12$). $5 \times 12 = 60$.

\begin{notaprofesor}[Referencia]
  Neugebauer, O. \textit{The Exact Sciences in Antiquity}, 2.ª ed.
  Dover, 1969. Es el texto clásico sobre astronomía y matemáticas del
  antiguo Oriente Próximo. Cap.\ 1 cubre el sistema sexagesimal.
\end{notaprofesor}

\subsection{Base 12 — Historia y argumentos}

La divisibilidad de 12 ha atraído a matemáticos durante siglos.
El escritor y matemático inglés Thomas Digges (siglo \textsc{xvi}) mencionó
ventajas del duodecimal. En el siglo \textsc{xviii}, durante el debate sobre
el sistema métrico en Francia, matemáticos como Jean-Charles de Borda
(quien promovió la división del círculo en 400 grados centesimales) y otros
reformadores discutieron activamente distintas bases.

La \textit{Dozenal Society of Great Britain} se fundó en 1959; su homóloga
americana, la \textit{Dozenal Society of America}, en 1944.
Sus publicaciones argumentan que la educación matemática temprana sería más
intuitiva en base 12, ya que los niños aprenderían fracciones sin que
aparezcan decimales infinitos para algo tan natural como ``un tercio''.

\begin{notaprofesor}[Para la discusión en clase]
  Una pregunta que genera debate productivo: \textit{``¿Cambiarías el sistema
  de numeración si pudieras? ¿Qué perderías?''}
  Respuestas interesantes: perderíamos compatibilidad con todo el software
  y hardware existente; ganaríamos fracciones más limpias. Es un buen
  momento para hablar de costos de migración en sistemas reales.
\end{notaprofesor}

% =============================================================================
\section{Preguntas frecuentes}
% =============================================================================

\begin{notaprofesor}[FAQ]

\textbf{``¿Por qué en C++ los literales hex llevan \texttt{0x}?''}\\
Es una convención del lenguaje C (1972) para distinguir hex de decimal sin
ambigüedad: \texttt{255} es decimal, \texttt{0xFF} es hex, \texttt{0377}
es octal. El compilador necesita saber la base para generar el valor
correcto en memoria.

\vspace{0.6em}
\textbf{``¿Qué tan grande puede ser un número en binario con $n$ bits?''}\\
Con $n$ bits puedes representar $2^n$ valores distintos: de 0 a $2^n - 1$.
Con 8 bits (1 byte): 0 a 255. Con 32 bits: 0 a $4{,}294{,}967{,}295$.
Con 64 bits: 0 a aproximadamente $1.8 \times 10^{19}$.

\vspace{0.6em}
\textbf{``¿Por qué se usa hexadecimal y no octal en los sistemas modernos?''}\\
Porque los procesadores modernos operan en múltiplos de 8 bits (bytes).
$8 = 2 \times 4$ bits, no $2 \times 3$ bits, así que 2 dígitos hex
representan exactamente 1 byte, mientras que octal no encaja limpiamente.
Octal era más común en los años 60–70 cuando había máquinas de 12 o 36 bits.

\vspace{0.6em}
\textbf{``¿Los negativos cómo se representan?''}\\
Eso es tema de otra sesión (complemento a dos), pero si alguien lo pregunta:
el bit más significativo se usa como signo; con $n$ bits se representan
valores de $-2^{n-1}$ a $2^{n-1}-1$.
\end{notaprofesor}

% =============================================================================
\section{Material de la sesión}
% =============================================================================

\begin{itemize}
  \item \textbf{Notas manuscritas originales}: \texttt{materialInicial/sesion03/}
  \item \textbf{Notas alumno}: \texttt{notas\_alumno/pdf/sesion03\_sistemas\_numeracion.pdf}
\end{itemize}

% =============================================================================
\section{Rúbrica de ejercicios}
% =============================================================================

\begin{center}
\begin{tabular}{lcc}
  \toprule
  \textbf{Criterio} & \textbf{Puntos} & \textbf{Evidencia} \\
  \midrule
  Ejercicios 1–4 correctos (1 punto c/u)          & 4 & resultado correcto \\
  Ejercicio 5 (octal → hex vía binario)            & 2 & proceso visible \\
  Ejercicio 6 (base 4 → hex vía binario)           & 2 & proceso visible \\
  Conversión directa sin pasar por decimal (5 y 6) & 2 & método binario intermedio \\
  \bottomrule
  \textbf{Total} & 10 & \\
\end{tabular}
\end{center}

\begin{notaprofesor}
  El bono de 2 puntos por conversión directa incentiva que el alumno
  interiorice la relación estructural entre las bases, en lugar de hacer
  dos conversiones separadas pasando por decimal.
  Si alguien llega al resultado correcto por decimal, es válido pero no
  obtiene esos 2 puntos.
\end{notaprofesor}

\end{document}
